\begin{thebibliography}{999}
% verified: CrossRef 10.1002/nme.995 2026-10-09
\bibitem{AbateValko04} J. Abate and P. P. Valk\'o. Multi-precision Laplace transform inversion. \emph{International Journal for Numerical Methods in Engineering} 60 (2004), 979--993.
% verified: arXiv 2608.13637 2026-10-11 (version 2 of 19 August 2026; title, authors and the constants 2/3 and 0.6725 checked in the abstract)
\bibitem{AF26} L. Alp\"oge and R. Furman. More than two thirds of the zeta zeros are simple and on the critical line. arXiv:2608.13637, 2026.
% verified: zbMATH 0456.60081 2026-09-27 (CrossRef 10.1137/1125045 dates the translated issue 25(2) January 1981; Russian original Teor. Veroyatnost. i Primenen. 25 (1980), 366)
\bibitem{Anulova80} S. V. Anulova. On Markov stopping times with a given distribution for a Wiener process. \emph{Theory of Probability and its Applications} 25 (1980), 362--366.
% verified: CrossRef 10.1090/s0025-5718-2010-02426-3 2026-09-27 (CrossRef splits the surname as 'de Reyna' and dates it online 2010; the AMS page and zbMATH 1228.11127 give Arias de Reyna, 80 (2011))
\bibitem{AriasDeReyna11} J. Arias de Reyna. High precision computation of Riemann's zeta function by the Riemann-Siegel formula, I. \emph{Mathematics of Computation} 80 (2011), 995--1009.
% verified: zbMATH 1291.46060 2026-10-10 (ALEA 10, no DOI in CrossRef); arXiv 1201.0311v2, text checked (Proposition 6, Definition 7, formula (7), Remark 8 and Theorem 9 in Section 3.2, as cited)
\bibitem{AHS13} O. Arizmendi, T. Hasebe and N. Sakuma. On the law of free subordinators. \emph{ALEA Latin American Journal of Probability and Mathematical Statistics} 10 (2013), 271--291.
% verified: CrossRef 10.1016/j.jcta.2017.11.012 2026-10-10
\bibitem{AP18} O. Arizmendi and D. Perales. Cumulants for finite free convolution. \emph{Journal of Combinatorial Theory, Series A} 155 (2018), 244--266.
% verified: CrossRef 10.1007/bf01404125 2026-10-09 (CrossRef dates the issue September 1971)
\bibitem{Armitage71} J. V. Armitage. Zeta functions with a zero at $s=\frac12$. \emph{Inventiones Mathematicae} 15 (1971), 199--205.
% verified: CrossRef 10.1214/13-ps221 2026-10-11; arXiv 1308.2729v3, text checked (Theorem 11.1 and the remark after it, Section 11.1)
\bibitem{AGK19} R. Arratia, L. Goldstein and F. Kochman. Size bias for one and all. \emph{Probability Surveys} 16 (2019), 1--61.
% verified: zbMATH 45.0719.02 2026-09-27 (JFM title reads 'Sur las'; title and pages 1979 to 1981 checked on the scan of C. R. 158, seance du 29 juin 1914)
\bibitem{Backlund14} R. J. Backlund. Sur les z\'eros de la fonction $\zeta(s)$ de Riemann. \emph{Comptes Rendus de l'Acad\'emie des Sciences} 158 (1914), 1979--1981.
% verified: CrossRef 10.1007/bf02422950 2026-09-26
\bibitem{Backlund} R. J. Backlund. \"Uber die Nullstellen der Riemannschen Zetafunktion. \emph{Acta Mathematica} 41 (1918), 345--375.
% verified: zbMATH 1097.11041 2026-09-26
\bibitem{BaezDuarte03} L. B\'aez-Duarte. A strengthening of the Nyman--Beurling criterion for the Riemann hypothesis. \emph{Atti della Accademia Nazionale dei Lincei, Rendiconti Lincei, Matematica e Applicazioni} 14 (2003), 5--11.
% verified: CrossRef 10.4064/aa-46-1-73-79 2026-10-09; zbMATH 0535.10045 (the review gives the exponent 7/22 on RH, after 9/28 of Montgomery and Vaughan and 8/25 of Graham)
\bibitem{BakerPintz85} R. C. Baker and J. Pintz. The distribution of square-free numbers. \emph{Acta Arithmetica} 46 (1985), 73--79.
% verified: CrossRef 10.1090/s0025-5718-07-01999-0 2026-10-11 (Math. Comp. 76, no. 260, 2045-2049); abstract (CrossRef and OpenAlex): zeros off the critical line of the Dirichlet series of Davenport and Heilbronn
\bibitem{BSO07} E. Balanzario and J. S\'anchez-Ortiz. Zeros of the Davenport--Heilbronn counterexample. \emph{Mathematics of Computation} 76 (2007), 2045--2049.
% verified: CrossRef 10.1016/j.laa.2018.03.012 2026-10-11; arXiv 1708.07729v3, text checked (Theorem 10, Section 4)
\bibitem{BS18} \'A. Baricz and F. \v{S}tampach. The Hurwitz-type theorem for the regular Coulomb wave function via Hankel determinants. \emph{Linear Algebra and its Applications} 548 (2018), 259--272.
% verified: zbMATH 1102.60013 2026-10-10 (not in CrossRef; the DOI suggested earlier, 10.3150/bj/1137421635, belongs to another Bernoulli paper)
\bibitem{BNMS06} O. E. Barndorff-Nielsen, M. Maejima and K. Sato. Some classes of multivariate infinitely divisible distributions admitting stochastic integral representations. \emph{Bernoulli} 12 (2006), 1--33.
% verified: zbMATH 1024.60022 2026-10-10 (not in CrossRef)
\bibitem{BNT02} O. E. Barndorff-Nielsen and S. Thorbj{\o}rnsen. Self-decomposability and L\'evy processes in free probability. \emph{Bernoulli} 8 (2002), 323--366.
% verified: CrossRef 10.1142/s0219025704001773 2026-10-10
\bibitem{BNT04} O. E. Barndorff-Nielsen and S. Thorbj{\o}rnsen. A connection between free and classical infinite divisibility. \emph{Infinite Dimensional Analysis, Quantum Probability and Related Topics} 7 (2004), 573--590.
% verified: CrossRef 10.2977/prims/31 2026-10-09 (the CrossRef abstract states the Sato Tate conjecture for non-CM holomorphic cuspidal newforms of weight greater than one)
\bibitem{BLGHT11} T. Barnet-Lamb, D. Geraghty, M. Harris and R. Taylor. A family of Calabi--Yau varieties and potential automorphy II. \emph{Publications of the Research Institute for Mathematical Sciences} 47 (2011), 29--98.
% verified: CrossRef 10.1090/mcom/3599 2026-10-10 (title, journal, volume 90, pages 1455 to 1482, issued 2021, authors Bennett, Martin, O'Bryant, Rechnitzer); arXiv 2005.02989 abstract checked for the bound used in Theorem thm:dbeta:minall(v): for conductor q > 1, T >= 5/7 and l = log(q(T+2)/(2 pi)) > 1.567, |N(T,chi) - ((T/pi) log(qT/(2 pi e)) - chi(-1)/4)| <= 0.22737 l + 2 log(1+l) - 0.5
\bibitem{BMOR21} M. A. Bennett, G. Martin, K. O'Bryant and A. Rechnitzer. Counting zeros of Dirichlet $L$-functions. \emph{Mathematics of Computation} 90 (2021), 1455--1482.

% verified: CrossRef 10.2307/121080 2026-10-10 (CrossRef gives the first page only; pages 1023 to 1060 from zbMATH 0945.46046; CrossRef lists P. Biane, the author of the appendix, as a third author)
\bibitem{BP99} H. Bercovici and V. Pata. Stable laws and domains of attraction in free probability theory (with an appendix by P. Biane). \emph{Annals of Mathematics} 149 (1999), 1023--1060.
% verified: CrossRef 10.1512/iumj.1993.42.42033 2026-10-10 (CrossRef gives no title and the first page only; title and pages 733 to 773 from zbMATH 0806.46070)
\bibitem{BV93} H. Bercovici and D. Voiculescu. Free convolution of measures with unbounded support. \emph{Indiana University Mathematics Journal} 42 (1993), 733--773.
% verified: CrossRef 10.1017/s1446788700008703 2026-10-11; text checked (the remark of [11, page 92] corrected on page 67); the first page header says 67--73, the last page is 74
\bibitem{Berg04} C. Berg. Correction to a paper by A. G. Pakes. \emph{Journal of the Australian Mathematical Society} 76 (2004), 67--74.
% verified: CrossRef 10.1007/BF00538793 2026-10-11 (Theorems 1.3 and 1.6 and Section 4 checked in a scan of the paper)
\bibitem{BergForst83} C. Berg and G. Forst. Multiply self-decomposable probability measures on $\mathbb R_+$ and $\mathbb Z_+$. \emph{Zeitschrift f\"ur Wahrscheinlichkeitstheorie und verwandte Gebiete} 62 (1983), 147--163.
% verified: CrossRef 10.1137/s0036144598347497 2026-10-09
\bibitem{BerryKeating99} M. V. Berry and J. P. Keating. The Riemann zeros and eigenvalue asymptotics. \emph{SIAM Review} 41 (1999), 236--266.
% verified: CrossRef 10.1073/pnas.41.5.312 2026-09-26
\bibitem{Beurling55} A. Beurling. A closure problem related to the Riemann zeta-function. \emph{Proceedings of the National Academy of Sciences} 41 (1955), 312--314.
% verified: CrossRef 10.1512/iumj.1997.46.1467 2026-10-09 (CrossRef and OpenAlex give no pages; pages 705 to 718 from the printed page numbers of the article PDF)
\bibitem{Biane97} P. Biane. On the free convolution with a semi-circular distribution. \emph{Indiana University Mathematics Journal} 46 (1997), 705--718.

% verified: CrossRef 10.1007/978-3-642-01763-6_7 2026-10-09 (pages 171 to 185; the Lecture Notes volume number is not in CrossRef or OpenAlex and is omitted); arXiv 0811.1490, text checked
\bibitem{Biane09} P. Biane. Matrix valued Brownian motion and a paper by P\'olya. In \emph{S\'eminaire de Probabilit\'es XLII}, Lecture Notes in Mathematics, Springer, 2009, 171--185.

% verified: CrossRef 10.1090/s0273-0979-01-00912-0 2026-09-26
\bibitem{BPY01} P. Biane, J. Pitman and M. Yor. Probability laws related to the Jacobi theta and Riemann zeta functions, and Brownian excursions. \emph{Bulletin of the American Mathematical Society} 38 (2001), 435--465.
% verified: CrossRef 10.1017/cbo9780511721434 2026-10-05 (Theorem 4.12.9 is de Bruijn's exponential Tauberian theorem, as cited by Aurzada and Dereich, Ann. Probab. 37 (2009), arXiv:0805.1330, reference [4]; the 1989 paperback keeps the numbering)
\bibitem{BGT87} N. H. Bingham, C. M. Goldie and J. L. Teugels. \emph{Regular Variation}. Encyclopedia of Mathematics and its Applications 27, Cambridge University Press, 1987.
% verified: CrossRef 10.1109/tau.1970.1162132 2026-09-27
\bibitem{Bluestein70} L. I. Bluestein. A linear filtering approach to the computation of discrete Fourier transform. \emph{IEEE Transactions on Audio and Electroacoustics} 18 (1970), 451--455.
% verified: CrossRef 10.1088/0305-4470/39/34/010 2026-09-27
\bibitem{BBLM06} E. Bogomolny, O. Bohigas, P. Leboeuf and A. G. Monastra. On the spacing distribution of the Riemann zeros: corrections to the asymptotic result. \emph{Journal of Physics A: Mathematical and General} 39 (2006), 10743--10754.
% verified: zbMATH 45.0716.02 2026-09-26
\bibitem{BL14} H. Bohr and E. Landau. Sur les z\'eros de la fonction $\zeta(s)$ de Riemann. \emph{Comptes Rendus de l'Acad\'emie des Sciences} 158 (1914), 106--110.
% verified: CrossRef 10.1006/jnth.1999.2392 2026-09-26
\bibitem{BL99} E. Bombieri and J. C. Lagarias. Complements to Li's criterion for the Riemann hypothesis. \emph{Journal of Number Theory} 77 (1999), 274--287.
% verified: zbMATH 1008.11034 2026-09-26
\bibitem{Bombieri} E. Bombieri. Remarks on Weil's quadratic functional in the theory of prime numbers, I. \emph{Rendiconti Lincei. Matematica e Applicazioni} (9) 11 (2000), 183--233.
% verified: CrossRef 10.1007/978-1-4612-2948-3 2026-09-26
\bibitem{Bondesson} L. Bondesson. \emph{Generalized Gamma Convolutions and Related Classes of Distributions and Densities}. Lecture Notes in Statistics 76, Springer, 1992.
% verified: CrossRef 10.1007/s10959-013-0523-y 2026-09-26
\bibitem{Bond2015} L. Bondesson. A class of probability distributions that is closed with respect to addition as well as multiplication of independent random variables. \emph{Journal of Theoretical Probability} 28 (2015), 1063--1081.
% verified: zbMATH 1366.60049 2026-09-26
\bibitem{boBosch} P. Bosch. HCM property and the half-Cauchy distribution. \emph{Probability and Mathematical Statistics} 35 (2015), 191--200.
% verified: CrossRef 10.1007/s11512-015-0216-0 2026-09-27
\bibitem{BoschSimon16} P. Bosch and T. Simon. A proof of Bondesson's conjecture on stable densities. \emph{Arkiv f\"or Matematik} 54 (2016), 31--38.
% verified: CrossRef 10.1214/ecp.v12-1244 2026-09-26; pages 73 to 80 from zbMATH 1129.60088, 2026-10-02
\bibitem{boBFY} P. Bourgade, T. Fujita and M. Yor. Euler's formulae for $\zeta(2n)$ and products of Cauchy variables. \emph{Electronic Communications in Probability} 12 (2007), 73--80.
% verified: CrossRef 10.1215/00127094-2008-046 2026-09-26 (pages 45 to 69 from the journal; CrossRef gives no pages)
\bibitem{BHNY08} P. Bourgade, C. P. Hughes, A. Nikeghbali and M. Yor. The characteristic polynomial of a random unitary matrix: a probabilistic approach. \emph{Duke Mathematical Journal} 145 (2008), 45--69.
% verified: zbMATH 1099.11511 2026-09-27
\bibitem{Bourgain00} J. Bourgain. On large values estimates for Dirichlet polynomials and the density hypothesis for the Riemann zeta function. \emph{International Mathematics Research Notices} 2000 (2000), 133--146.
% verified: CrossRef 10.1137/1028165 2026-10-11 (title 'A Conjectured Definite Integral (A. H. Nuttall)'; the problem number is not in the record)
\bibitem{Bouwkamp86} C. J. Bouwkamp. A conjectured definite integral (A. H. Nuttall), solution. \emph{SIAM Review} 28 (1986), 568--569.
% verified: CrossRef 10.1090/S0025-5718-1979-0537983-2 2026-09-27
\bibitem{Brent79} R. P. Brent. On the zeros of the Riemann zeta function in the critical strip. \emph{Mathematics of Computation} 33 (1979), 1361--1372.
% verified: CrossRef 10.1090/mcom/3643 2026-09-30
\bibitem{BKLNW21} S. Broadbent, H. Kadiri, A. Lumley, N. Ng and K. Wilk. Sharper bounds for the Chebyshev function $\theta(x)$. \emph{Mathematics of Computation} 90 (2021), 2281--2315.
% verified: CrossRef 10.1111/1467-9965.00035 2026-10-09
\bibitem{BGK97} M. Broadie, P. Glasserman and S. Kou. A continuity correction for discrete barrier options. \emph{Mathematical Finance} 7 (1997), 325--349.
% verified: CrossRef 10.1007/pl00004786 2026-09-26
\bibitem{BCKV} D. Bump, K.-K. Choi, P. Kurlberg and J. Vaaler. A local Riemann hypothesis, I. \emph{Mathematische Zeitschrift} 233 (2000), 1--18.
% verified: CrossRef 10.1090/mcom/3264 2026-09-30 (published online 26 October 2017; volume 87, number 312, is the 2018 volume)
\bibitem{Buthe18} J. B\"uthe. An analytic method for bounding $\psi(x)$. \emph{Mathematics of Computation} 87 (2018), 1991--2009.
% verified: arXiv 2512.11405 2026-09-26
\bibitem{Candelpergher} B. Candelpergher. A new expansion of the Riemann zeta function. arXiv:2512.11405, 2025.
% verified: CrossRef 10.1007/s00229-003-0414-0 2026-09-26
\bibitem{Cardon2004} D. A. Cardon. Sums of exponential functions having only real zeros. \emph{Manuscripta Mathematica} 113 (2004), 307--317.
% verified: CrossRef 10.1090/s0002-9939-04-07677-4 2026-09-25
\bibitem{Cardon2005} D. A. Cardon. Fourier transforms having only real zeros. \emph{Proceedings of the American Mathematical Society} 133 (2005), 1349--1356.
% verified: CrossRef 10.1214/10-aap714 2026-09-27 (pages 1663 to 1693 from zbMATH 1244.60078 and arXiv:1112.5305; CrossRef gives no pages)
\bibitem{CCCS11} X. Chen, L. Cheng, J. Chadam and D. Saunders. Existence and uniqueness of solutions to the inverse boundary crossing problem for diffusions. \emph{Annals of Applied Probability} 21 (2011), 1663--1693.
% verified: CrossRef 10.1090/s0002-9947-1962-0143257-8 2026-10-12 (CrossRef gives the first page only; pages 434 to 450 as in Pitman and Yor 2003 [17]; the AMS page returned 403, so the paper itself was not read)
\bibitem{CiesielskiTaylor62} Z. Ciesielski and S. J. Taylor. First passage times and sojourn times for Brownian motion in space and the exact Hausdorff measure of the sample path. \emph{Transactions of the American Mathematical Society} 103 (1962), 434--450.
% verified: CrossRef 10.1090/S0025-5718-2013-02681-6 2026-10-09 (abstract from OpenAlex)
\bibitem{CC13} M. W. Coffey and G. Csordas. On the log-concavity of a Jacobi theta function. \emph{Mathematics of Computation} 82 (2013), 2265--2272.
% verified: zbMATH 0966.11039 2026-09-26
\bibitem{CL00} J. B. Conrey and X.-J. Li. A note on some positivity conditions related to zeta and $L$-functions. \emph{International Mathematics Research Notices} 2000 (2000), 929--940.
% verified: CrossRef 10.1112/s0024611504015175 2026-09-26
\bibitem{CFKRS05} J. B. Conrey, D. W. Farmer, J. P. Keating, M. O. Rubinstein and N. C. Snaith. Integral moments of $L$-functions. \emph{Proceedings of the London Mathematical Society} 91 (2005), 33--104.
% verified: CrossRef 10.1007/bf02124750 2026-10-11; full text checked (author's copy: the series (3.1), the expansion (4.19), the boundary curve (4.4) of the range of the principal branch in Section 4)
\bibitem{Corless96} R. M. Corless, G. H. Gonnet, D. E. G. Hare, D. J. Jeffrey and D. E. Knuth. On the Lambert $W$ function. \emph{Advances in Computational Mathematics} 5 (1996), 329--359.
% verified: CrossRef 10.1007/bf01203390 2026-09-27
\bibitem{Cramer19} H. Cram\'er. Studien \"uber die Nullstellen der Riemannschen Zetafunktion. \emph{Mathematische Zeitschrift} 4 (1919), 104--130.
% verified: CrossRef 10.1007/bf01482072 2026-09-26
\bibitem{Cramer} H. Cram\'er. Ein Mittelwertsatz in der Primzahltheorie. \emph{Mathematische Zeitschrift} 12 (1922), 147--153.
% verified: CrossRef 10.2307/1971315 2026-10-01; pages 405 to 431 from zbMATH 0625.30036, 2026-10-02
\bibitem{CCS87} T. Craven, G. Csordas and W. Smith. The zeros of derivatives of entire functions and the P\'olya-Wiman conjecture. \emph{Annals of Mathematics} 125 (1987), 405--431.
% verified: CrossRef 10.1090/s0002-9947-1986-0846596-4 2026-09-27
\bibitem{CNV86} G. Csordas, T. S. Norfolk and R. S. Varga. The Riemann hypothesis and the Tur\'an inequalities. \emph{Transactions of the American Mathematical Society} 296 (1986), 521--541.
% verified: zbMATH 0807.11059 2026-09-26
\bibitem{COSV93} G. Csordas, A. M. Odlyzko, W. Smith and R. S. Varga. A new Lehmer pair of zeros and a new lower bound for the de Bruijn-Newman constant $\Lambda$. \emph{Electronic Transactions on Numerical Analysis} 1 (1993), 104--111.
% verified: CrossRef 10.1007/bf01205170 2026-09-26
\bibitem{CSV94} G. Csordas, W. Smith and R. S. Varga. Lehmer pairs of zeros, the de Bruijn-Newman constant $\Lambda$, and the Riemann hypothesis. \emph{Constructive Approximation} 10 (1994), 107--129.
% verified: zbMATH 1002.11001 2026-09-26
\bibitem{Davenport} H. Davenport. \emph{Multiplicative Number Theory}, third edition, revised by H. L. Montgomery. Graduate Texts in Mathematics 74, Springer, 2000.
% verified: CrossRef 10.1112/jlms/s1-11.3.181 (part I, s1-11 no. 3, 181-185, July 1936) and 10.1112/jlms/s1-11.4.307 (part II, s1-11 no. 4, 307-312, October 1936), 2026-10-11; text not seen, so the part that contains the example of period 5 is unchecked and both are cited
\bibitem{DH36} H. Davenport and H. Heilbronn. On the zeros of certain Dirichlet series I, II. \emph{Journal of the London Mathematical Society} 11 (1936), 181--185 and 307--312.
% verified: CrossRef 10.1215/s0012-7094-50-01720-0 2026-09-26; pages 197 to 226 from zbMATH 0038.23302, 2026-10-02
\bibitem{deBruijn} N. G. de Bruijn. The roots of trigonometric integrals. \emph{Duke Mathematical Journal} 17 (1950), 197--226.
% verified: CrossRef 10.4064/cm-53-2-333-335 2026-09-27
\bibitem{Delange87} H. Delange. Une remarque sur la d\'eriv\'ee logarithmique de la fonction z\^eta de Riemann. \emph{Colloquium Mathematicum} 53 (1987), 333--335.
% verified: zbMATH 0287.14001 2026-09-26
\bibitem{Deligne74} P. Deligne. La conjecture de Weil. I. \emph{Publications Math\'ematiques de l'IH\'ES} 43 (1974), 273--307.
% verified: CrossRef 10.1007/978-1-4613-8643-8 2026-09-26
\bibitem{Devroye86} L. Devroye. \emph{Non-Uniform Random Variate Generation}. Springer, 1986.
% verified: CrossRef 10.4064/aa200603-23-7 2026-10-08
\bibitem{Dobner21} A. Dobner. A proof of Newman's conjecture for the extended Selberg class. \emph{Acta Arithmetica} 201 (2021), 29--62.
% verified: zbMATH 0327.60029 2026-09-26
\bibitem{DymMcKean76} H. Dym and H. P. McKean. \emph{Gaussian Processes, Function Theory, and the Inverse Spectral Problem}. Probability and Mathematical Statistics 31, Academic Press, 1976.
% verified: CrossRef 10.1063/1.1703862 2026-10-09
\bibitem{Dyson62} F. J. Dyson. A Brownian-motion model for the eigenvalues of a random matrix. \emph{Journal of Mathematical Physics} 3 (1962), 1191--1198.
% verified: zbMATH 0315.10035 2026-09-26
\bibitem{Edwards} H. M. Edwards. \emph{Riemann's Zeta Function}. Academic Press, 1974.
% verified: CrossRef 10.1214/16-aap1172 2026-09-27 (pages 3154 to 3177 from zbMATH 1351.60110 and Project Euclid; CrossRef gives no pages)
\bibitem{EkstromJanson16} E. Ekstr\"om and S. Janson. The inverse first-passage problem and optimal stopping. \emph{Annals of Applied Probability} 26 (2016), 3154--3177.
% verified: CrossRef 10.1007/bf01608495 2026-10-01
\bibitem{EMN76} R. S. Ellis, J. L. Monroe and C. M. Newman. The GHS and other correlation inequalities for a class of even ferromagnets. \emph{Communications in Mathematical Physics} 46 (1976), 167--182.
% verified: CrossRef 10.1007/bf00535504 2026-09-26
\bibitem{EGV79} P. Embrechts, C. M. Goldie and N. Veraverbeke. Subexponentiality and infinite divisibility. \emph{Zeitschrift f\"ur Wahrscheinlichkeitstheorie und Verwandte Gebiete} 49 (1979), 335--347.
% verified: zbMATH 0219.60003 2026-09-26
\bibitem{Feller71} W. Feller. \emph{An Introduction to Probability Theory and Its Applications}, volume II, second edition. Wiley, 1971.
% verified: arXiv 1604.07705 2026-10-01
\bibitem{Fourati16} S. Fourati. On some analytical properties of stable densities. arXiv:1604.07705, 2016.
% verified: CrossRef 10.1103/physrevd.7.2850 2026-10-09
\bibitem{Fulling73} S. A. Fulling. Nonuniqueness of canonical field quantization in Riemannian space-time. \emph{Physical Review D} 7 (1973), 2850--2862.
% verified: CrossRef 10.53846/goediss-5113 2026-09-25
\bibitem{Gabcke79} W. Gabcke. \emph{Neue Herleitung und explizite Restabsch\"atzung der Riemann-Siegel-Formel}. Dissertation, Georg-August-Universit\"at G\"ottingen, 1979.
% verified: CrossRef 10.1017/s0305004100042225 2026-10-09
\bibitem{Goldie67} C. Goldie. A class of infinitely divisible random variables. \emph{Proceedings of the Cambridge Philosophical Society} 63 (1967), 1141--1143.

% verified: zbMATH 1297.11101 2026-09-26
\bibitem{GonekMontgomery13} S. M. Gonek and H. L. Montgomery. Zeros of a family of approximations of the Riemann zeta-function. \emph{International Mathematics Research Notices} 2013 (2013), 4712--4733.
% verified: publisher URL http://numbers.computation.free.fr/Constants/Miscellaneous/zetazeros1e13-1e24.pdf 2026-09-27
\bibitem{Gourdon04} X. Gourdon. The $10^{13}$ first zeros of the Riemann Zeta function, and zeros computation at very large height. Manuscript, 2004, \url{http://numbers.computation.free.fr/Constants/Miscellaneous/zetazeros1e13-1e24.pdf}, accessed 27 September 2026.
% verified: CrossRef 10.1007/bf02421310 2026-09-27
\bibitem{Gram03} J.-P. Gram. Note sur les z\'eros de la fonction $\zeta(s)$ de Riemann. \emph{Acta Mathematica} 27 (1903), 289--304.
% verified: CrossRef 10.1073/pnas.1902572116 2026-10-01
\bibitem{GORZ19} M. Griffin, K. Ono, L. Rolen and D. Zagier. Jensen polynomials for the Riemann zeta function and other sequences. \emph{Proceedings of the National Academy of Sciences} 116 (2019), 11103--11110.
% verified: CrossRef 10.1063/1.1665211 2026-10-01
\bibitem{GHS70} R. B. Griffiths, C. A. Hurst and S. Sherman. Concavity of magnetization of an Ising ferromagnet in a positive external field. \emph{Journal of Mathematical Physics} 11 (1970), 790--795.
% verified: CrossRef 10.1007/bf02465533 2026-10-12 (English translation of Liet. Mat. Rink. 39 (1999), 40 to 51; the publisher's page returned a client challenge, so the paper itself was not read)
\bibitem{Grigelionis99} B. Grigelionis. Processes of Meixner type. \emph{Lithuanian Mathematical Journal} 39 (1999), 33--41.
% verified: CrossRef 10.1515/crll.1914.144.114 2026-10-11; text not checked (cited through BS18, Theorem 10 (i), reference [8])
\bibitem{Grommer14} J. Grommer. Ganze transzendente Funktionen mit lauter reellen Nullstellen. \emph{Journal f\"ur die reine und angewandte Mathematik} 144 (1914), 114--166.
% verified: CrossRef 10.4007/annals.2026.203.2.6 2026-09-27 (pages 623 to 675 from the journal; CrossRef gives no pages)
\bibitem{GM26} L. Guth and J. Maynard. New large value estimates for Dirichlet polynomials. \emph{Annals of Mathematics} 203 (2026), 623--675.
% verified: zbMATH 46.0498.01 2026-09-26
\bibitem{HL16} G. H. Hardy and J. E. Littlewood. Contributions to the theory of the Riemann zeta-function and the theory of the distribution of primes. \emph{Acta Mathematica} 41 (1916), 119--196.
% verified: zbMATH 1159.11001 2026-09-26
\bibitem{boHardyWright} G. H. Hardy and E. M. Wright. \emph{An Introduction to the Theory of Numbers}, sixth edition, revised by D. R. Heath-Brown and J. H. Silverman. Oxford University Press, 2008.
% verified: CrossRef 10.1214/aop/1176996606 2026-10-09 (CrossRef gives no pages; pages 593 to 607 from zbMATH)
\bibitem{HartmanWatson74} P. Hartman and G. S. Watson. ``Normal'' distribution functions on spheres and the modified Bessel functions. \emph{Annals of Probability} 2 (1974), 593--607.
% verified: CrossRef 10.1016/j.jnt.2021.06.032 2026-10-11 (title, journal, volume 235, pages 219 to 241, issued June 2022, authors Hasanalizade, Shen, Wong); arXiv 2107.06506v1 read for (1.7), (2.7), Corollaries 1.2 and 1.4 and footnote 2 (the error in Trudgian 2014)
\bibitem{HSW22} E. Hasanalizade, Q. Shen and P.-J. Wong. Counting zeros of the Riemann zeta function. \emph{Journal of Number Theory} 235 (2022), 219--241.
% verified: CrossRef 10.1007/s10959-015-0595-y 2026-10-10
\bibitem{HT16} T. Hasebe and S. Thorbj{\o}rnsen. Unimodality of the freely selfdecomposable probability laws. \emph{Journal of Theoretical Probability} 29 (2016), 922--940.
% verified: CrossRef 10.1007/bf01565437 2026-10-09
\bibitem{Hecke36} E. Hecke. \"Uber die Bestimmung Dirichletscher Reihen durch ihre Funktionalgleichung. \emph{Mathematische Annalen} 112 (1936), 664--699.

% verified: CrossRef 10.1215/s0012-7094-97-08601-4 2026-09-26; pages 1 to 37 from zbMATH 0887.46008, 2026-10-02
\bibitem{HLS97} H. Hedenmalm, P. Lindqvist and K. Seip. A Hilbert space of Dirichlet series and systems of dilated functions in $L^2(0,1)$. \emph{Duke Mathematical Journal} 86 (1997), 1--37.
% verified: CrossRef 10.1515/crll.1856.52.39 2026-10-11; text not checked (cited as BS18 does, reference [9], for the theorem on polynomials attributed to Hermite)
\bibitem{Hermite1856} C. Hermite. Extrait d'une lettre de Mr. Ch. Hermite de Paris \`a Mr. Borchardt de Berlin sur le nombre des racines d'une \'equation alg\'ebrique comprises entre des limites donn\'ees. \emph{Journal f\"ur die reine und angewandte Mathematik} 52 (1856), 39--51.
% verified: CrossRef 10.1016/0550-3213(85)90418-3 2026-10-09
\bibitem{tHooft85} G. 't~Hooft. On the quantum structure of a black hole. \emph{Nuclear Physics B} 256 (1985), 727--745.
% verified: CrossRef 10.1090/mcom/3275 2026-09-27
\bibitem{Hurst18} G. Hurst. Computations of the Mertens function and improved bounds on the Mertens conjecture. \emph{Mathematics of Computation} 87 (2018), 1013--1028.
% verified: CrossRef 10.1090/s0002-9947-1925-1501297-5 2026-09-27
\bibitem{Hutchinson25} J. I. Hutchinson. On the roots of the Riemann zeta function. \emph{Transactions of the American Mathematical Society} 27 (1925), 49--60.
% verified: zbMATH 0241.10026 2026-09-27 (year 1972 as in zbMATH and MR 292774; CrossRef 10.1007/bf01418933 dates the issue June 1971, before the paper's receipt on 4 November 1971)
\bibitem{Huxley72} M. N. Huxley. On the difference between consecutive primes. \emph{Inventiones Mathematicae} 15 (1972), 164--170.
% verified: CrossRef 10.1007/978-0-8176-4532-8_5 2026-10-09 (CrossRef gives no year for the chapter; the volume, CrossRef 10.1007/978-0-8176-4532-8, edited by V. Ginzburg, is dated 2006; the series number 253 is not in CrossRef and is taken from the note)
\bibitem{Ihara06} Y. Ihara. On the Euler--Kronecker constants of global fields and primes with small norms. In \emph{Algebraic Geometry and Number Theory}, V. Ginzburg, ed., Progress in Mathematics 253, Birkh\"auser, 2006, 407--451.
% verified: CrossRef 10.1214/009117904000000225 2026-10-11 (title, authors, volume 32, issue 2); pages 1356 to 1369 from the journal reference of arXiv math/0205316 (CrossRef and OpenAlex give none)
\bibitem{IJS04} A. M. Iksanov, Z. J. Jurek and B. M. Schreiber. A new factorization property of the selfdecomposable probability measures. \emph{Annals of Probability} 32 (2004), 1356--1369.
% verified: CrossRef 10.1112/plms/s2-27.1.273 2026-09-26
\bibitem{Ingham26} A. E. Ingham. Mean-value theorems in the theory of the Riemann zeta-function. \emph{Proceedings of the London Mathematical Society} (2) 27 (1928), 273--300.
% verified: zbMATH 0025.02704 2026-09-26
\bibitem{Ingham40} A. E. Ingham. On the estimation of $N(\sigma,T)$. \emph{Quarterly Journal of Mathematics} 11 (1940), 291--292.
% verified: CrossRef 10.2307/2371685 2026-10-01; pages 313 to 319 from zbMATH 0063.02974, 2026-10-02
\bibitem{Ingham42} A. E. Ingham. On two conjectures in the theory of numbers. \emph{American Journal of Mathematics} 64 (1942), 313--319.
% verified: arXiv 1412.1220 2026-09-26
\bibitem{Inoue} H. Inoue. Expansion of Dirichlet $L$-function on the critical line in Meixner-Pollaczek polynomials. arXiv:1412.1220, 2014.
% verified: zbMATH 0127.09503 2026-09-26
\bibitem{ItoMcKean65} K. It\^o and H. P. McKean. \emph{Diffusion Processes and Their Sample Paths}. Springer, 1965.
% verified: zbMATH 1059.11001 2026-09-27
\bibitem{IwaniecKowalski04} H. Iwaniec and E. Kowalski. \emph{Analytic Number Theory}. American Mathematical Society Colloquium Publications 53, American Mathematical Society, 2004.
% verified: CrossRef 10.1214/07-ps118 2026-09-26; pages 346 to 415 from zbMATH 1189.60035, 2026-10-02
\bibitem{JRY} L. F. James, B. Roynette and M. Yor. Generalized gamma convolutions, Dirichlet means, Thorin measures, with explicit examples. \emph{Probability Surveys} 5 (2008), 346--415.
% verified: CrossRef 10.1109/tc.2017.2690633 2026-09-26
\bibitem{Arb} F. Johansson. Arb: efficient arbitrary-precision midpoint-radius interval arithmetic. \emph{IEEE Transactions on Computers} 66 (2017), 1281--1292.
% verified: CrossRef 10.1016/j.jmaa.2023.127460 2026-09-30
\bibitem{JY23} D. R. Johnston and A. Yang. Some explicit estimates for the error term in the prime number theorem. \emph{Journal of Mathematical Analysis and Applications} 527 (2023), 127460.
% verified: zbMATH 1112.58038 2026-10-09 (no DOI and no arXiv record; journal, volume, year and pages as in reference [10] of Biane09)
\bibitem{JonesOConnell06} L. Jones and N. O'Connell. Weyl chambers, symmetric spaces and number variance saturation. \emph{ALEA Latin American Journal of Probability and Mathematical Statistics} 2 (2006), 91--118.
% verified: CrossRef 10.1007/bf00538800 2026-09-26
\bibitem{JurekVervaat83} Z. J. Jurek and W. Vervaat. An integral representation for selfdecomposable Banach space valued random variables. \emph{Zeitschrift f\"ur Wahrscheinlichkeitstheorie und Verwandte Gebiete} 62 (1983), 247--262.
% verified: CrossRef 10.2307/2035888 2026-09-26; pages 1155 to 1157 from zbMATH 0136.33001, 2026-10-02
\bibitem{Kac65} M. Kac. A remark on Wiener's Tauberian theorem. \emph{Proceedings of the American Mathematical Society} 16 (1965), 1155--1157.
% verified: arXiv 1011.5940v1 2026-10-11 (27 November 2010, no journal reference; OpenAlex lists no published version)
\bibitem{KJ10} G. A. Kalugin and D. J. Jeffrey. Unimodal sequences show Lambert $W$ is Bernstein. arXiv:1011.5940, 2010.
% verified: arXiv 1103.5640v1 2026-10-11; text checked (relations (10) and (11), Theorems 3.3, 6.2 and 6.5, Corollary 3.4 and the remark after it)
\bibitem{KJC11} G. A. Kalugin, D. J. Jeffrey and R. M. Corless. Stieltjes, Poisson and other integral representations for functions of Lambert $W$. Preprint, arXiv:1103.5640, 2011; published as \cite{KJCB12,KJC12}.
% verified: CrossRef 10.1080/10652469.2011.613830 2026-10-11 (issue 8); abstract checked
\bibitem{KJCB12} G. A. Kalugin, D. J. Jeffrey, R. M. Corless and P. B. Borwein. Stieltjes and other integral representations for functions of Lambert $W$. \emph{Integral Transforms and Special Functions} 23 (2012), 581--593.
% verified: CrossRef 10.1080/10652469.2011.640327 2026-10-11 (issue 11, online 14 December 2011)
\bibitem{KJC12} G. A. Kalugin, D. J. Jeffrey and R. M. Corless. Bernstein, Pick, Poisson and related integral expressions for Lambert $W$. \emph{Integral Transforms and Special Functions} 23 (2012), 817--829.
% verified: CrossRef 10.1007/s002200000261 2026-09-26
\bibitem{KS00} J. P. Keating and N. C. Snaith. Random matrix theory and $\zeta(1/2+it)$. \emph{Communications in Mathematical Physics} 214 (2000), 57--89.
% verified: CrossRef 10.1090/s0025-5718-1992-1122072-5 2026-09-26; pages 765 to 773 from zbMATH 0767.11039; CrossRef gives 765 to 765, 2026-10-02
\bibitem{Keiper92} J. B. Keiper. Power series expansions of Riemann's $\xi$ function. \emph{Mathematics of Computation} 58 (1992), 765--773.
% verified: CrossRef 10.1214/aop/1176995427 2026-10-01 (CrossRef gives no pages; pages 760 to 770 from zbMATH 0402.60080)
\bibitem{Kent78} J. T. Kent. Some probabilistic properties of Bessel functions. \emph{Annals of Probability} 6 (1978), 760--770.
% verified: CrossRef 10.1007/bf00538895 2026-09-26
\bibitem{Kent80} J. T. Kent. Eigenvalue expansions for diffusion hitting times. \emph{Zeitschrift f\"ur Wahrscheinlichkeitstheorie und Verwandte Gebiete} 52 (1980), 309--319.
% verified: CrossRef 10.1214/aop/1176993924 2026-10-09 (CrossRef gives no pages; pages 207 to 219 from zbMATH 0483.60075; statement checked in the Project Euclid abstract)
\bibitem{Kent82} J. T. Kent. The spectral decomposition of a diffusion hitting time. \emph{Annals of Probability} 10 (1982), 207--219.

% verified: CrossRef 10.1090/s0025-5718-1983-0717706-5 2026-09-29
\bibitem{Kershaw83} D. Kershaw. Some extensions of W. Gautschi's inequalities for the gamma function. \emph{Mathematics of Computation} 41 (1983), 607--611.
% verified: CrossRef 10.1016/j.aim.2009.04.003 2026-09-29
\bibitem{KKL09} H. Ki, Y.-O. Kim and J. Lee. On the de Bruijn--Newman constant. \emph{Advances in Mathematics} 222 (2009), 281--306.
% verified: CrossRef 10.1007/978-3-642-05014-5 2026-09-25
\bibitem{KLS} R. Koekoek, P. A. Lesky and R. F. Swarttouw. \emph{Hypergeometric Orthogonal Polynomials and Their $q$-Analogues}. Springer, 2010.
% verified: CrossRef 10.5802/aif.3600 2026-09-27
\bibitem{KPS24} T. Konstantopoulos, P. Patie and R. Sarkar. A new class of solutions to the van Dantzig problem, the Lee--Yang property, and the Riemann hypothesis. \emph{Annales de l'Institut Fourier} 74 (2024), 377--421.
% verified: CrossRef 10.2307/2034651 2026-09-26; pages 353 to 355 from zbMATH 0134.11401, 2026-10-02
\bibitem{Korevaar65} J. Korevaar. Distribution proof of Wiener's Tauberian theorem. \emph{Proceedings of the American Mathematical Society} 16 (1965), 353--355.
% verified: CrossRef 10.1214/aop/1176988867 2026-10-09 (pages 442 to 452 from zbMATH 0796.60017)
\bibitem{Kristiansen94} G. K. Kristiansen. A proof of Steutel's conjecture. \emph{Annals of Probability} 22 (1994), 442--452.

% verified: CrossRef 10.1006/jath.1999.3316 2026-10-10
\bibitem{KVA99} A. B. J. Kuijlaars and W. Van Assche. The asymptotic zero distribution of orthogonal polynomials with varying recurrence coefficients. \emph{Journal of Approximation Theory} 99 (1999), 167--197.
% verified: CrossRef 10.1080/10652460701450773 2026-09-26
\bibitem{Kuz2007} A. Kuznetsov. Integral representations for the Dirichlet $L$-functions and their expansions in Meixner--Pollaczek polynomials and rising factorials. \emph{Integral Transforms and Special Functions} 18 (2007), 827--835.
% verified: CrossRef 10.4153/cmb-2008-055-0 2026-09-26
\bibitem{Kuz2008} A. Kuznetsov. Expansion of the Riemann $\Xi$ function in Meixner--Pollaczek polynomials. \emph{Canadian Mathematical Bulletin} 51 (2008), 561--569.
% verified: CrossRef 10.4064/aa-89-3-217-234 2026-09-26
\bibitem{Lagarias} J. C. Lagarias. On a positivity property of the Riemann $\xi$-function. \emph{Acta Arithmetica} 89 (1999), 217--234.
% verified: CrossRef 10.4064/aa120-2-4 2026-10-01 (statement checked in the arXiv abstract, math/0601653)
\bibitem{Lagarias05} J. C. Lagarias. Zero spacing distributions for differenced $L$-functions. \emph{Acta Arithmetica} 120 (2005), 159--184.
% verified: CrossRef 10.1111/j.1467-842x.1979.tb01143.x 2026-09-26
\bibitem{LVJ} C. D. Lai and D. Vere-Jones. Odd man out: the Meixner hypergeometric distribution. \emph{Australian Journal of Statistics} 21 (1979), 256--265.
% verified: CrossRef 10.1090/s0002-9947-1958-0094863-x 2026-09-26
\bibitem{boLamperti} J. Lamperti. An occupation time theorem for a class of stochastic processes. \emph{Transactions of the American Mathematical Society} 88 (1958), 380--387.
% verified: arXiv 2609.02882 2026-10-11 (version 2 of 8 September 2026; title, author and the constant 67.25 per cent checked in the abstract)
\bibitem{Lam26} Y. Lamzouri. A new proof that more than $2/3$ of the zeros of the Riemann zeta function are simple and on the critical line. arXiv:2609.02882, 2026.
% verified: CrossRef 10.1007/978-1-4612-0853-2 2026-10-09 (second edition, Graduate Texts in Mathematics 110)
\bibitem{Lang94} S. Lang. \emph{Algebraic Number Theory}, second edition. Graduate Texts in Mathematics 110, Springer, 1994.
% verified: zbMATH 0860.65028 2026-09-26
\bibitem{LawsonHanson74} C. L. Lawson and R. J. Hanson. \emph{Solving Least Squares Problems}. Prentice Hall, 1974.
% verified: zbMATH 0138.29101 2026-10-09 (the paper is in the combined issue numbers 159 to 170 of the series, pages 41 to 48; the note cites it as number 170 under a Serbian title)
\bibitem{Lazarevic66} I. Lazarevi\'c. Certaines in\'egalit\'es o\`u interviennent les fonctions hyperboliques. \emph{Univerzitet u Beogradu, Publikacije Elektrotehni\v{c}kog Fakulteta, Serija Matematika i Fizika} 159--170 (1966), 41--48.
% verified: CrossRef 10.1103/physrev.87.410 2026-09-26
\bibitem{LeeYang} T. D. Lee and C. N. Yang. Statistical theory of equations of state and phase transitions. II. Lattice gas and Ising model. \emph{Physical Review} 87 (1952), 410--419.
% verified: CrossRef 10.1112/plms/s3-20.2.303 2026-09-26
\bibitem{Lehman70} R. S. Lehman. On the distribution of zeros of the Riemann zeta-function. \emph{Proceedings of the London Mathematical Society} (3) 20 (1970), 303--320.
% verified: CrossRef 10.1007/bf02401102 2026-09-27
\bibitem{Lehmer56} D. H. Lehmer. On the roots of the Riemann zeta-function. \emph{Acta Mathematica} 95 (1956), 291--298.
% verified: CrossRef 10.1112/s0025579300001753 2026-09-27
\bibitem{Lehmer56b} D. H. Lehmer. Extended computation of the Riemann zeta-function. \emph{Mathematika} 3 (1956), 102--108.
% verified: CrossRef 10.1006/jnth.1997.2137 2026-09-26
\bibitem{Li97} X.-J. Li. The positivity of a sequence of numbers and the Riemann hypothesis. \emph{Journal of Number Theory} 65 (1997), 325--333.
% verified: CrossRef 10.1007/bf01213009 2026-10-01
\bibitem{LiebSokal81} E. H. Lieb and A. D. Sokal. A general Lee-Yang theorem for one-component and multicomponent ferromagnets. \emph{Communications in Mathematical Physics} 80 (1981), 153--179.
% verified: CrossRef 10.1090/tran/7249 2026-09-26
\bibitem{LPS18} A. Lindner, L. Pan and K. Sato. On quasi-infinitely divisible distributions. \emph{Transactions of the American Mathematical Society} 370 (2018), 8483--8520.
% verified: zbMATH 0121.35502 2026-09-27
\bibitem{Linnik64} Yu. V. Linnik. \emph{Decomposition of Probability Distributions}. Oliver and Boyd, 1964.
% verified: publisher URL https://www.lmfdb.org/citation 2026-09-27 (zeros of L-function 2-1-1.1-c11-0-0, modular form 1.12.a.a, checked on https://www.lmfdb.org/L/ModularForm/GL2/Q/holomorphic/1/12/a/a/; LMFDB Release 1.2.1)
\bibitem{LMFDB} The LMFDB Collaboration. \emph{The $L$-functions and modular forms database}, \url{https://www.lmfdb.org}, accessed 27 September 2026.
% verified: CrossRef 10.1137/1113008 2026-09-29
\bibitem{Lukacs68} E. Lukacs. Contributions to a problem of D. van Dantzig. \emph{Theory of Probability and its Applications} 13 (1968), 116--127.
% verified: zbMATH 0207.07201 2026-09-26
\bibitem{boMandelbrojt} S. Mandelbrojt. \emph{S\'eries de Dirichlet. Principes et m\'ethodes}. Gauthier-Villars, 1969.
% verified: CrossRef 10.1070/SM1967v001n04ABEH001994 2026-10-10
\bibitem{MP67} V. A. Mar\v{c}enko and L. A. Pastur. Distribution of eigenvalues for some sets of random matrices. \emph{Mathematics of the USSR, Sbornik} 1 (1967), 457--483.
% verified: CrossRef 10.4007/annals.2015.182.1.8 2026-10-10 (CrossRef gives no volume; volume 182 from the DOI)
\bibitem{MSS15} A. W. Marcus, D. A. Spielman and N. Srivastava. Interlacing families II: mixed characteristic polynomials and the Kadison Singer problem. \emph{Annals of Mathematics} 182 (2015), 327--350.
% verified: CrossRef 10.1007/s00440-021-01105-w 2026-10-09; arXiv 1504.00350, text checked (the form (p +_d q)(x) = p^(D) q^(D) x^d after Definition 1.1, the Hermite polynomials H_d = e^{-D^2/2} x^d in Section 1.2, and the remark there that limits of the finite convolutions give the free ones, citing Marcus 2015 and Arizmendi and Perales 2018)
\bibitem{MSS22} A. W. Marcus, D. A. Spielman and N. Srivastava. Finite free convolutions of polynomials. \emph{Probability Theory and Related Fields} 182 (2022), 807--848.

% verified: CrossRef 10.1007/bf02504648 2026-09-27
\bibitem{Mase75} S. Mase. Decomposition of infinitely divisible characteristic functions with absolutely continuous Poisson spectral measure. \emph{Annals of the Institute of Statistical Mathematics} 27 (1975), 289--298.
% verified: Stack Exchange API 2.3, question 2063166 on math.stackexchange.com, 2026-10-09 (title; asked by the user Desiderius Severus on 18 December 2016; accepted answer 2063794 by Peter Humphries, with the factor 1/k in its Euler factor; second answer 2064069 by the user reuns; comment by Marco Cantarini)
\bibitem{MSE2063166} Math StackExchange question 2063166. Dirichlet series associated to squared M\"obius. December 2016, \url{https://math.stackexchange.com/q/2063166}, accessed 9 October 2026.
% verified: CrossRef 10.1016/0001-8708(83)90035-x 2026-09-26; pages 95 to 125 from OpenAlex, 2026-10-02
\bibitem{MR} N. Metropolis and G.-C. Rota. Witt vectors and the algebra of necklaces. \emph{Advances in Mathematics} 50 (1983), 95--125.
% verified: CrossRef 10.1090/pspum/024/9944 2026-09-27
\bibitem{Montgomery73} H. L. Montgomery. The pair correlation of zeros of the zeta function. In \emph{Analytic Number Theory}, H. G. Diamond, ed., Proceedings of Symposia in Pure Mathematics 24, American Mathematical Society, 1973, 181--193.
% verified: PyPI mpmath-1.3.0 2026-09-26
\bibitem{mpmath} The mpmath development team. \emph{mpmath: a Python library for arbitrary-precision floating-point arithmetic}, version 1.3.0, 2023.
% verified: CrossRef 10.3150/13-BEJ581 2026-10-01 (pages 604 to 617 from the journal line of the arXiv reprint 1504.03438; CrossRef gives no pages)
\bibitem{Nakamura15} T. Nakamura. A complete Riemann zeta distribution and the Riemann hypothesis. \emph{Bernoulli} 21 (2015), 604--617.
% verified: CrossRef 10.1002/cpa.3160270203 2026-09-27
\bibitem{Newman74} C. M. Newman. Zeros of the partition function for generalized Ising systems. \emph{Communications on Pure and Applied Mathematics} 27 (1974), 143--159.
% verified: CrossRef 10.1007/bf01608542 2026-10-01
\bibitem{Newman75} C. M. Newman. Inequalities for Ising models and field theories which obey the Lee-Yang theorem. \emph{Communications in Mathematical Physics} 41 (1975), 1--9.
% verified: CrossRef 10.1090/s0002-9939-1976-0434982-5 2026-09-25
\bibitem{Newman1976} C. M. Newman. Fourier transforms with only real zeros. \emph{Proceedings of the American Mathematical Society} 61 (1976), 245--251.
% verified: CrossRef 10.1007/bf01888165 2026-10-01
\bibitem{Newman91} C. M. Newman. The GHS inequality and the Riemann hypothesis. \emph{Constructive Approximation} 7 (1991), 389--399.
% verified: CrossRef 10.1090/bull/1668 2026-10-09 (arXiv 1901.06596, text checked: Theorem 10 there is the characterisation of Newman1976 used in the draft)
\bibitem{NewmanWu20} C. M. Newman and W. Wu. Constants of de Bruijn--Newman type in analytic number theory and statistical physics. \emph{Bulletin of the American Mathematical Society} 57 (2020), 595--614.
% verified: CrossRef 10.1017/CBO9780511735127 2026-10-10; zbMATH 1133.60003 (series and volume number)
\bibitem{NS06} A. Nica and R. Speicher. \emph{Lectures on the Combinatorics of Free Probability}. London Mathematical Society Lecture Note Series 335, Cambridge University Press, 2006.
% verified: CrossRef 10.1214/ecp.v14-1488 2026-09-26 (pages 396 to 411 from the journal; CrossRef gives no pages)
\bibitem{NY09} A. Nikeghbali and M. Yor. The Barnes $G$ function and its relations with sums and products of generalized gamma convolution variables. \emph{Electronic Communications in Probability} 14 (2009), 396--411.
% verified: CrossRef 10.1137/s0040585x9797571x 2026-10-12 (English translation, issued 1997; Russian original Teor. Veroyatnost. i Primenen. 41 (1996), 810 to 826, CrossRef 10.4213/tvp3203, read at mathnet.ru: Theorem 2, formula (7))
\bibitem{Novikov96} A. A. Novikov. Martingales, Tauberian theorem, and strategies of gambling. \emph{Theory of Probability and its Applications} 41 (1997), 716--729.
% verified: CrossRef 10.1137/1027144 2026-10-11 (title 'A Conjectured Definite Integral'; the problem number is not in the record)
\bibitem{Nuttall85} A. H. Nuttall. A conjectured definite integral, problem. \emph{SIAM Review} 27 (1985), 573.
% verified: zbMATH 0037.35401 2026-09-26
\bibitem{Nyman50} B. Nyman. \emph{On the One-Dimensional Translation Group and Semi-Group in Certain Function Spaces}. Thesis, University of Uppsala, 1950.
% verified: CrossRef 10.1090/s0025-5718-1987-0866115-0 2026-09-25
\bibitem{Odlyzko87} A. M. Odlyzko. On the distribution of spacings between zeros of the zeta function. \emph{Mathematics of Computation} 48 (1987), 273--308.
% verified: CrossRef 10.1090/s0002-9947-1988-0961614-2 2026-09-26
\bibitem{OS88} A. M. Odlyzko and A. Sch\"onhage. Fast algorithms for multiple evaluations of the Riemann zeta function. \emph{Transactions of the American Mathematical Society} 309 (1988), 797--809.
% verified: URL https://www-users.cse.umn.edu/~odlyzko/unpublished/zeta.10to20.1992.pdf 2026-09-27
\bibitem{Odlyzko92} A. M. Odlyzko. \emph{The $10^{20}$-th zero of the Riemann zeta function and 175 million of its neighbors}. Unpublished manuscript, 1992, \url{https://www-users.cse.umn.edu/~odlyzko/unpublished/zeta.10to20.1992.pdf}, accessed 27 September 2026.
% verified: CrossRef 10.1023/a:1016677511798 2026-09-29 (the bound -2.63e-9 is computed after (3.14) in Section 3; the abstract and (1.12) state -2.7e-9)
\bibitem{Odlyzko00} A. M. Odlyzko. An improved bound for the de Bruijn--Newman constant. \emph{Numerical Algorithms} 25 (2000), 293--303.
% verified: CrossRef 10.1090/conm/290/04578 2026-09-27
\bibitem{Odlyzko01} A. M. Odlyzko. The $10^{22}$-nd zero of the Riemann zeta function. In \emph{Dynamical, Spectral, and Arithmetic Zeta Functions}, M. L. Lapidus and M. van Frankenhuysen, eds., Contemporary Mathematics 290, American Mathematical Society, 2001, 139--144.
% verified: URL https://www-users.cse.umn.edu/~odlyzko/zeta_tables/ 2026-09-27
\bibitem{OdlyzkoTables} A. M. Odlyzko. \emph{Tables of zeros of the Riemann zeta function}. Online, \url{https://www-users.cse.umn.edu/~odlyzko/zeta_tables/}, accessed 27 September 2026.
% verified: publisher URL https://dlmf.nist.gov/help/cite 2026-09-27
\bibitem{DLMF} F. W. J. Olver, A. B. Olde Daalhuis, D. W. Lozier, B. I. Schneider, R. F. Boisvert, C. W. Clark, B. R. Miller, B. V. Saunders, H. S. Cohl and M. A. McClain, eds. \emph{NIST Digital Library of Mathematical Functions}. \url{https://dlmf.nist.gov/}, Release 1.2.8 of 2026-09-15.

\bibitem{OAI26a} OpenAI. The quasi-Riemann hypothesis: a zero-free half-plane $\re(s)>7/8$. Preprint dated 30 September 2026, released 6 October 2026 with a formalisation in Lean, \url{https://github.com/openai/math} (family 003), accessed 7 October 2026.

\bibitem{OAI26b} OpenAI. The quasi-Riemann hypothesis. Preprint dated 5 October 2026, released 6 October 2026, \url{https://github.com/openai/math} (family 003), accessed 7 October 2026.
% verified: CrossRef 10.1006/jmaa.1996.0056 2026-10-11
\bibitem{Pakes96} A. G. Pakes. Length biasing and laws equivalent to the log-normal. \emph{Journal of Mathematical Analysis and Applications} 197 (1996), 825--854.
% verified: CrossRef 10.1017/s1446788700002731 2026-10-11; abstract checked
\bibitem{Pakes01} A. G. Pakes. Remarks on converse Carleman and Krein criteria for the classical moment problem. \emph{Journal of the Australian Mathematical Society} 71 (2001), 81--104.
% verified: CrossRef 10.1016/j.jmaa.2006.03.035 2026-10-11
\bibitem{Pakes07} A. G. Pakes. Structure of Stieltjes classes of moment-equivalent probability laws. \emph{Journal of Mathematical Analysis and Applications} 326 (2007), 1268--1290.
% verified: CrossRef 10.1016/j.jmaa.2011.01.066 2026-10-11; text not checked (content as described in Pakes26, Section 2)
\bibitem{Pakes11} A. G. Pakes. Lambert's $W$, infinite divisibility and Poisson mixtures. \emph{Journal of Mathematical Analysis and Applications} 378 (2011), 480--492.
% verified: CrossRef 10.1155/2014/703697 2026-10-11 (article 703697, 9 pages); abstract checked
\bibitem{Pakes14} A. G. Pakes. On the preservation of infinite divisibility under length-biasing. \emph{Journal of Probability} 2014 (2014), 703697.
% verified: CrossRef 10.1016/j.spl.2018.03.015 2026-10-11; text not checked (content as described in Pakes26, Section 2)
\bibitem{Pakes18} A. G. Pakes. The Lambert $W$ function, Nuttall's integral, and the Lambert law. \emph{Statistics and Probability Letters} 139 (2018), 53--60.
% verified: CrossRef 10.1007/s10959-019-00886-0 2026-10-11 (online 2019); text not checked (content as described in Pakes26, Section 2)
\bibitem{Pakes20} A. G. Pakes. Self-decomposable laws from continuous branching processes. \emph{Journal of Theoretical Probability} 33 (2020), 361--395.
% verified: CrossRef 10.3390/axioms15080574 2026-10-11; text checked (Sections 1, 2 and 6, Theorem 1, equation (53))
\bibitem{Pakes26} A. G. Pakes. Self-decomposability and the Lambert law. \emph{Axioms} 15 (2026), 574.
% verified: CrossRef 10.1111/1467-842x.00158 2026-10-11
\bibitem{PHW01} A. G. Pakes, W.-L. Hung and J.-W. Wu. Criteria for the unique determination of probability distributions by moments. \emph{Australian and New Zealand Journal of Statistics} 43 (2001), 101--111.
% verified: CrossRef 10.1111/j.1467-842x.1992.tb01363.x 2026-10-11; abstract checked
\bibitem{PK92} A. G. Pakes and R. Khattree. Length-biasing, characterizations of laws and the moment problem. \emph{Australian Journal of Statistics} 34 (1992), 307--322.
% verified: CrossRef 10.1007/bf02926159 2026-10-11; text not checked (cited by Pakes26 for the length bias characterisation)
\bibitem{PSF96} A. G. Pakes, T. Sapatinas and E. B. Fosam. Characterizations, length-biasing, and infinite divisibility. \emph{Statistical Papers} 37 (1996), 53--69.
% verified: CrossRef 10.1214/ECP.v13-1413 2026-10-10
\bibitem{PAS08} V. P\'erez-Abreu and N. Sakuma. Free generalized gamma convolutions. \emph{Electronic Communications in Probability} 13 (2008), 526--539.
% verified: CrossRef 10.1216/jiea/1181074930 2026-09-26; pages 397 to 423 from zbMATH 1044.60076, 2026-10-02
\bibitem{Peskir02} G. Peskir. On integral equations arising in the first-passage problem for Brownian motion. \emph{Journal of Integral Equations and Applications} 14 (2002), 397--423.
% verified: CrossRef 10.1007/bfb0088732 2026-10-09 (pages 285 to 370; the volume number 851 from zbMATH 0469.60076)
\bibitem{PitmanYor81} J. Pitman and M. Yor. Bessel processes and infinitely divisible laws. In \emph{Stochastic Integrals}, Lecture Notes in Mathematics 851, Springer, 1981, 285--370.
% verified: CrossRef 10.4153/cjm-2003-014-x 2026-10-12; PDF read at Cambridge Core (formulas (1), (4), (67) and Table 5 as cited)
\bibitem{PitmanYor03} J. Pitman and M. Yor. Infinitely divisible laws associated with hyperbolic functions. \emph{Canadian Journal of Mathematics} 55 (2003), 292--330.
% verified: zbMATH 1345.11064 2026-09-26
\bibitem{Platt16} D. J. Platt. Numerical computations concerning the GRH. \emph{Mathematics of Computation} 85 (2016), 3009--3027.
% verified: CrossRef 10.1112/blms.12460 2026-09-26
\bibitem{PT21} D. J. Platt and T. S. Trudgian. The Riemann hypothesis is true up to $3\cdot10^{12}$. \emph{Bulletin of the London Mathematical Society} 53 (2021), 792--797.
% verified: CrossRef 10.1111/rssb.12042 2026-09-27
\bibitem{PSW14} N. G. Polson, J. G. Scott and J. Windle. The Bayesian bridge. \emph{Journal of the Royal Statistical Society, Series B (Statistical Methodology)} 76 (2014), 713; sections and theorems are cited by the numbering of arXiv:1109.2279v2.
% verified: arXiv 1708.02653v24 2026-09-27
\bibitem{P17} N. G. Polson. On Hilbert's 8th problem. arXiv:1708.02653v24, 2026.
% verified: arXiv 1708.02653v5 2026-09-27
\bibitem{P17v5} N. G. Polson. On Hilbert's 8th problem. arXiv:1708.02653v5, 2017.
% verified: publisher URL https://www.amazon.com/dp/B0HKPL1X1S 2026-09-27 (Independently published, 22 September 2026, 488 pages, ISBN 9798176411898; title and subtitle from the title page of the interior PDF, whose copyright page reads "First edition, 2026" and prints no publisher or ISBN; pinpoints follow this edition, not the draft of about 195 pages)
\bibitem{P26} N. G. Polson. \emph{On Hilbert's 8th Problem: the reciprocal Xi clock, Thorin measures, van Dantzig pairs, zero free approximations, hitting times and the reciprocal zeta function}, first edition. Independently published, 2026.
% verified: manuscript N/A 2026-10-09 (title, sole author and date from the source file dirichlet_beta_thorin.tex; no arXiv record)
\bibitem{P26beta} N. G. Polson. The Dirichlet beta function in Thorin form. Manuscript, October 2026.
% verified: manuscript N/A 2026-10-09 (title and sole author from the source file biane_polya_thorin.tex, whose title page is dated September 2026; the files were revised on 9 October 2026; no arXiv record)
\bibitem{P26bia} N. G. Polson. Biane on P\'olya: hitting times, Thorin measures and a Wick rotation. Manuscript, October 2026.
% verified: manuscript N/A 2026-09-27 (title, sole author and date from the source file eta_paper.tex; no arXiv record)
\bibitem{P26eta} N. G. Polson. Thorin measures of the Dedekind eta function. Manuscript, September 2026.
% verified: manuscript N/A 2026-10-08 (title, sole author and date from the source file max_exponential_incomplete_gamma.tex; no arXiv record)
\bibitem{P26max} N. G. Polson. Maxima of exponentials, incomplete gamma continuations, and the Gamma(2) structure of $L$ functions. Manuscript, October 2026.
% verified: manuscript N/A 2026-09-27 (title, byline and date from reciprocal-zeta/paper/reciprocal_zeta_thorin.tex)
\bibitem{P26rz} N. G. Polson. Thorin type representations of the reciprocal zeta function on the critical line. Manuscript, September 2026, now Part IX of the first edition of \emph{On Hilbert's 8th Problem} (2026).
% verified: manuscript N/A 2026-10-09 (title, sole author and date from the source file squarefree_xi_thorin.tex; no arXiv record)
\bibitem{P26sqf} N. G. Polson. The squarefree Dirichlet series $\zeta(s)/\zeta(2s)$: a summary of Math StackExchange 2063166, P\'olya densities, Gamma(2) smoothing and Thorin measures of $\Xi$. Manuscript, October 2026.
% verified: manuscript N/A 2026-10-09 (title, sole author and date from the source file thorin_identities.tex; no arXiv record)
\bibitem{P26thor} N. G. Polson. Gamma, reciprocal Gamma, $\eta^3$ and $\Delta$: Thorin forms and the Riemann $\Xi$ function. Manuscript, October 2026.
% verified: manuscript N/A 2026-10-09 (title, sole author and date from the source files wigner_free_thorin_v3.tex and wigner_free_thorin_v7.tex, the latter dated October 2026 (version 7); no arXiv record)
\bibitem{P26wig} N. G. Polson. Wigner's semicircle, free probability and the Thorin framework for the reciprocal Xi function. Manuscript, version 7, October 2026.
% verified: manuscript N/A 2026-10-11 (title, sole author and date from the source file pakes_lambert_thorin.tex, dated October 2026; the chapter version pakes_chapter.tex is titled 'Pakes, the Lambert Law and the Thorin Condition'; no arXiv record)
\bibitem{P26pak} N. G. Polson. Pakes, the Lambert law and the Thorin condition for the Riemann Xi function. Manuscript, October 2026.
% verified: manuscript N/A 2026-09-27 (authors' source thorin_xi_computation.tex, title page dated September 2026; not on arXiv)
\bibitem{PS26} N. G. Polson and V. Sokolov. Computing the Thorin measure of the reciprocal Xi function: a P\'olya route and a total variation comparison with the Riemann atoms. Manuscript, September 2026.
% verified: CrossRef 10.1007/bf01199419 2026-10-01
\bibitem{Polya18} G. P\'olya. \"Uber die Nullstellen gewisser ganzer Funktionen. \emph{Mathematische Zeitschrift} 2 (1918), 352--383.
% verified: zbMATH 49.0228.01 2026-09-26
\bibitem{boPolya23} G. P\'olya. \"Uber die Existenz unendlich vieler singul\"arer Punkte auf der Konvergenzgeraden gewisser Dirichletscher Reihen. \emph{Sitzungsberichte der Preussischen Akademie der Wissenschaften, Physikalisch-mathematische Klasse} (1923), 45--50.
% verified: CrossRef 10.1007/bf02565336 2026-09-26
\bibitem{Polya} G. P\'olya. Bemerkung \"uber die Integraldarstellung der Riemannschen $\xi$-Funktion. \emph{Acta Mathematica} 48 (1926), 305--317.
% verified: CrossRef 10.1007/bf01180553 2026-09-26
\bibitem{boPolya} G. P\'olya. Untersuchungen \"uber L\"ucken und Singularit\"aten von Potenzreihen. \emph{Mathematische Zeitschrift} 29 (1929), 549--640.
% verified: zbMATH 0319.01021 2026-09-27 (Vol. II is Mathematicians of Our Time 8; Kac's comments at pp. 424 to 426 per Newman, PAMS 61 (1976), 245, and arXiv:0809.0401)
\bibitem{Kac74} G. P\'olya. \emph{Collected Papers, Vol.~II: Location of Zeros}, edited by R. P. Boas. Mathematicians of Our Time 8, MIT Press, 1974.
% verified: CrossRef 10.1515/crll.1914.144.89 2026-10-10
\bibitem{PS14} G. P\'olya and J. Schur. \"Uber zwei Arten von Faktorenfolgen in der Theorie der algebraischen Gleichungen. \emph{Journal f\"ur die reine und angewandte Mathematik} 144 (1914), 89--113.
% verified: CrossRef 10.1007/s40687-019-0193-1 2026-09-26
\bibitem{Polymath} D. H. J. Polymath. Effective approximation of heat flow evolution of the Riemann $\xi$ function, and a new upper bound for the de Bruijn--Newman constant. \emph{Research in the Mathematical Sciences} 6 (2019), paper 31.
% verified: CrossRef 10.1007/s40687-019-0199-8 2026-10-11 (volume 7, issue 1, article number 2; CrossRef and OpenAlex date it online 6 December 2019, and the journal reference of arXiv 1802.10521 gives volume 7 (2020))
\bibitem{PRZZ20} K. Pratt, N. Robles, A. Zaharescu and D. Zeindler. More than five-twelfths of the zeros of $\zeta$ are on the critical line. \emph{Research in the Mathematical Sciences} 7 (2020), article 2.
% verified: zbMATH 0895.30001 2026-09-27
\bibitem{Remmert98} R. Remmert. \emph{Classical Topics in Complex Function Theory}, translated by L. Kay. Graduate Texts in Mathematics 172, Springer, 1998.
% verified: CrossRef 10.1007/bf02418544 2026-09-26
\bibitem{Riesz16} M. Riesz. Sur l'hypoth\`ese de Riemann. \emph{Acta Mathematica} 40 (1916), 185--190.
% verified: CrossRef 10.2307/2308012 2026-09-28; pages 26 to 29 from MathSciNet MR Lookup, 2026-10-02
\bibitem{Robbins55} H. Robbins. A remark on Stirling's formula. \emph{American Mathematical Monthly} 62 (1955), 26--29.
% verified: CrossRef 10.1214/aop/1176990230 2026-09-27 (pages 1689 to 1731 from zbMATH 0739.60074; CrossRef gives no pages)
\bibitem{Roberts91} G. O. Roberts. Asymptotic approximations for Brownian motion boundary hitting times. \emph{Annals of Probability} 19 (1991), 1689--1731.
% verified: CrossRef 10.1214/aoap/1177004773 2026-09-26; pages 446 to 460 from zbMATH 0836.60087, 2026-10-02
\bibitem{RS95} G. O. Roberts and C. F. Shortland. The hazard rate tangent approximation for boundary hitting times. \emph{Annals of Applied Probability} 5 (1995), 446--460.
% verified: CrossRef 10.1017/fmp.2020.6 2026-09-26
\bibitem{RodgersTao} B. Rodgers and T. Tao. The de Bruijn--Newman constant is non-negative. \emph{Forum of Mathematics, Pi} 8 (2020), e6.
% verified: CrossRef 10.4064/aa200515-10-3 2026-09-26
\bibitem{Romik} D. Romik. Orthogonal polynomial expansions for the Riemann xi function in the Hermite, Meixner--Pollaczek, and continuous Hahn bases. \emph{Acta Arithmetica} 200 (2021), 259--329.
% verified: CrossRef 10.1215/ijm/1255631807 2026-09-28 (pages 64 to 94 from zbMATH 0122.05001; CrossRef gives no pages)
\bibitem{RS62} J. B. Rosser and L. Schoenfeld. Approximate formulas for some functions of prime numbers. \emph{Illinois Journal of Mathematics} 6 (1962), 64--94.
% verified: zbMATH 0191.44504 2026-09-27 (editor and publisher from Open Library and from Brent, Math. Comp. 33 (1979), reference 23); pages 70 to 76 from zbMATH 0191.44504, 2026-10-02
\bibitem{RYS69} J. B. Rosser, J. M. Yohe and L. Schoenfeld. Rigorous computation and the zeros of the Riemann zeta-function. In \emph{Information Processing 68: Proceedings of IFIP Congress 1968}, volume 1, A. J. H. Morrell, ed., North-Holland, 1969, 70--76.
% verified: zbMATH 0792.11034 2026-09-26; pages 415 to 440 from zbMATH 0792.11034, which adds the supplement S17 to S23, 2026-10-02
\bibitem{Rumely93} R. Rumely. Numerical computations concerning the ERH. \emph{Mathematics of Computation} 61 (1993), 415--440.
% verified: zbMATH 0050.08003 2026-09-26
\bibitem{Salem53} R. Salem. Sur une proposition \'equivalente \`a l'hypoth\`ese de Riemann. \emph{Comptes Rendus de l'Acad\'emie des Sciences} 236 (1953), 1127--1128.
% verified: CrossRef 10.1016/j.anihpb.2006.09.007 2026-10-12; arXiv math/0512440, text checked (formula (2.5) in Section 2)
\bibitem{SalminenVallois07} P. Salminen and P. Vallois. On maximum increase and decrease of Brownian motion. \emph{Annales de l'Institut Henri Poincar\'e, Probabilit\'es et Statistiques} 43 (2007), 655--676.
% verified: CrossRef 10.1016/0047-259X(80)90014-7 2026-10-11 (Theorem 3.5, p. 219, checked in the publisher PDF by the triage, check/C_literature.md)
\bibitem{Sato80} K. Sato. Class $L$ of multivariate distributions and its subclasses. \emph{Journal of Multivariate Analysis} 10 (1980), 207--232.
% verified: zbMATH 0973.60001 2026-09-26
\bibitem{Sato99} K. Sato. \emph{L\'evy Processes and Infinitely Divisible Distributions}. Cambridge Studies in Advanced Mathematics 68, Cambridge University Press, 1999.
% verified: CrossRef 10.1214/14-ss107 2026-10-09 (pages 45 to 114 from OpenAlex; CrossRef gives no pages; preservation under convolution is Proposition 3.5 of arXiv:1404.5886)
\bibitem{SW14} A. Saumard and J. A. Wellner. Log-concavity and strong log-concavity: a review. \emph{Statistics Surveys} 8 (2014), 45--114.
% verified: CrossRef 10.1515/9783110269338 2026-09-26
\bibitem{boSSV} R. L. Schilling, R. Song and Z. Vondra\v{c}ek. \emph{Bernstein Functions: Theory and Applications}, second edition. Walter de Gruyter, 2012.
% verified: CrossRef 10.1007/bf02790092 2026-09-27 (accent in Polya from zbMATH 0045.37602 and the original first page; CrossRef omits it)
\bibitem{Schoenberg51} I. J. Schoenberg. On P\'olya frequency functions. I. The totally positive functions and their Laplace transforms. \emph{Journal d'Analyse Math\'ematique} 1 (1951), 331--374.
% verified: CrossRef 10.1080/15326349808807475 2026-10-12 (pages 335 to 349 as in Pitman and Yor 2003 [64] and Schoutens 2001 [17]; the publisher's page returned 403, so the paper itself was not read)
\bibitem{SchoutensTeugels98} W. Schoutens and J. L. Teugels. L\'evy processes, polynomials and martingales. \emph{Communications in Statistics. Stochastic Models} 14 (1998), 335--349.
% verified: zbMATH 0061.08402 2026-09-26
\bibitem{Selberg46} A. Selberg. Contributions to the theory of the Riemann zeta-function. \emph{Archiv for Mathematik og Naturvidenskab} 48 (1946), 89--155.
% verified: zbMATH 0004.10501 2026-09-26; pages 45 to 80 from zbMATH 0004.10501, 2026-10-02
\bibitem{Siegel32} C. L. Siegel. \"Uber Riemanns Nachla{\ss} zur analytischen Zahlentheorie. \emph{Quellen und Studien zur Geschichte der Mathematik, Astronomie und Physik, Abteilung B: Studien} 2 (1932), 45--80.
% verified: CrossRef 10.1088/1751-8113/47/32/325204 2026-10-09
\bibitem{Sierra14} G. Sierra. The Riemann zeros as energy levels of a Dirac fermion in a potential built from the prime numbers in Rindler spacetime. \emph{Journal of Physics A: Mathematical and Theoretical} 47 (2014), 325204.
% verified: CrossRef 10.1007/bf01645626 2026-10-01
\bibitem{SimonGriffiths73} B. Simon and R. B. Griffiths. The $(\phi^4)_2$ field theory as a classical Ising model. \emph{Communications in Mathematical Physics} 33 (1973), 145--164.
% verified: CrossRef 10.1007/s10998-010-1037-3 2026-09-26
\bibitem{SD10} J. Sondow and C. Dumitrescu. A monotonicity property of Riemann's xi function and a reformulation of the Riemann hypothesis. \emph{Periodica Mathematica Hungarica} 60 (2010), 37--40.
% verified: zbMATH 0226.10040 2026-09-26
\bibitem{Spira67} R. Spira. Zeros of approximate functional approximations. \emph{Mathematics of Computation} 21 (1967), 41--48.
% verified: CrossRef 10.1090/s0025-5718-1994-1254148-8 2026-10-11 (Math. Comp. 63, no. 208, first page 747; CrossRef gives no last page); 747--748 as in Ivic, arXiv:math/0311162, ref. [58]
\bibitem{Spira94} R. Spira. Some zeros of the Titchmarsh counterexample. \emph{Mathematics of Computation} 63 (1994), 747--748.
% verified: CrossRef 10.1214/aoms/1177698806 2026-10-09
\bibitem{Steutel67} F. W. Steutel. Note on the infinite divisibility of exponential mixtures. \emph{Annals of Mathematical Statistics} 38 (1967), 1303--1305.
% verified: CrossRef 10.1016/0304-4149(73)90008-2 2026-10-11; text not checked (content as described in AGK19, Section 11.1)
\bibitem{Steutel73} F. W. Steutel. Some recent results in infinite divisibility. \emph{Stochastic Processes and their Applications} 1 (1973), 125--143.

% verified: CrossRef 10.1111/j.1467-9574.1989.tb01246.x 2026-10-09
\bibitem{SteutelThiemann89} F. W. Steutel and J. G. F. Thiemann. On the independence of integer and fractional parts. \emph{Statistica Neerlandica} 43 (1989), 53--59.
% verified: CrossRef 10.1201/9780203014127 2026-09-26
\bibitem{SvH} F. W. Steutel and K. van Harn. \emph{Infinite Divisibility of Probability Distributions on the Real Line}. Marcel Dekker, 2004.
% verified: CrossRef 10.1093/imamat/23.1.97 2026-09-26
\bibitem{Talbot79} A. Talbot. The accurate numerical inversion of Laplace transforms. \emph{Journal of the Institute of Mathematics and its Applications} 23 (1979), 97--120.
% verified: zbMATH 1492.11154 2026-09-27 (title as printed on p. 305; zbMATH adds a comma after fields)
\bibitem{Tate67} J. T. Tate. Fourier analysis in number fields and Hecke's zeta-functions. In \emph{Algebraic Number Theory}, J. W. S. Cassels and A. Fr\"ohlich, eds., Academic Press, 1967, 305--347.
% verified: CrossRef 10.1080/03461238.1977.10405623 2026-09-27 (CrossRef and the publisher misspell divisibility)
\bibitem{Thorin77a} O. Thorin. On the infinite divisibility of the Pareto distribution. \emph{Scandinavian Actuarial Journal} (1977), 31--40.
% verified: CrossRef 10.1080/03461238.1977.10405635 2026-09-26
\bibitem{Thorin} O. Thorin. On the infinite divisibility of the lognormal distribution. \emph{Scandinavian Actuarial Journal} (1977), 121--148.
% verified: CrossRef 10.1080/03461238.1978.10432021 2026-09-27
\bibitem{Thorin78} O. Thorin. An extension of the notion of a generalized $\Gamma$-convolution. \emph{Scandinavian Actuarial Journal} (1978), 141--149.
% verified: CrossRef 10.1098/rspa.1936.0192 2026-09-27
\bibitem{Titchmarsh36} E. C. Titchmarsh. The zeros of the Riemann zeta-function. \emph{Proceedings of the Royal Society of London, Series A} 157 (1936), 261--263.
% verified: Biane09 reference [17] (arXiv 0811.1490, text checked: Oxford, at the Clarendon Press, 1946); Part II followed in 1958 and the second edition of the first part, as Part I, in 1962 (CrossRef 10.1063/1.3062231, 10.1063/1.3058324, reviews); no DOI for the book
\bibitem{TitchmarshEigen46} E. C. Titchmarsh. \emph{Eigenfunction Expansions Associated with Second-Order Differential Equations}. Clarendon Press, Oxford, 1946.
% verified: zbMATH 0601.10026 2026-09-26
\bibitem{Titchmarsh} E. C. Titchmarsh. \emph{The Theory of the Riemann Zeta-Function}, second edition, revised by D. R. Heath-Brown. Oxford University Press, 1986.
% verified: CrossRef 10.1137/130932132 2026-09-26
\bibitem{TW14} L. N. Trefethen and J. A. C. Weideman. The exponentially convergent trapezoidal rule. \emph{SIAM Review} 56 (2014), 385--458.
% verified: CrossRef 10.1090/s0025-5718-2011-02470-1 2026-09-26
\bibitem{Trudgian11} T. S. Trudgian. Improvements to Turing's method. \emph{Mathematics of Computation} 80 (2011), 2259--2279.
% verified: CrossRef 10.4064/aa148-3-2 2026-09-27
\bibitem{Trudgian11b} T. S. Trudgian. On the success and failure of Gram's Law and the Rosser Rule. \emph{Acta Arithmetica} 148 (2011), 225--256.
% verified: CrossRef 10.1016/j.jnt.2013.07.017 2026-09-26
\bibitem{Trudgian} T. S. Trudgian. An improved upper bound for the argument of the Riemann zeta-function on the critical line II. \emph{Journal of Number Theory} 134 (2014), 280--292.
% verified: CrossRef 10.4064/aa-46-4-369-395 2026-09-26
\bibitem{Tsang} K.-M. Tsang. Some $\Omega$-theorems for the Riemann zeta-function. \emph{Acta Arithmetica} 46 (1986), 369--395.
% verified: CrossRef 10.1007/bf01459119 2026-10-11; text not checked (cited through BS18, Theorem 10 (ii), reference [29])
\bibitem{Chebotarev28} N. Tschebotareff. \"Uber die Realit\"at von Nullstellen ganzer transzendenter Funktionen. \emph{Mathematische Annalen} 99 (1928), 660--686.
% verified: CrossRef 10.1112/plms/s3-3.1.99 2026-09-26
\bibitem{Turing53} A. M. Turing. Some calculations of the Riemann zeta-function. \emph{Proceedings of the London Mathematical Society} (3) 3 (1953), 99--117.
% verified: CrossRef 10.1016/b978-0-12-426653-7.50021-1 2026-10-11 (OpenAlex W27910133; pages 225 to 237)
\bibitem{Urbanik73} K. Urbanik. Limit laws for sequences of normed sums satisfying some stability conditions. In \emph{Multivariate Analysis III}, P. R. Krishnaiah, ed., Academic Press, 1973, 225--237.
% verified: CrossRef 10.1090/s0025-5718-1986-0829637-3 2026-09-27
\bibitem{LRW86} J. van de Lune, H. J. J. te Riele and D. T. Winter. On the zeros of the Riemann zeta function in the critical strip. IV. \emph{Mathematics of Computation} 46 (1986), 667--681.
% verified: CrossRef 10.1016/0022-1236(86)90062-5 2026-10-10
\bibitem{Voi86} D. Voiculescu. Addition of certain non-commuting random variables. \emph{Journal of Functional Analysis} 66 (1986), 323--346.
% verified: zbMATH 0049.03205 2026-09-26; pages 252 to 265 from zbMATH 0049.03205, 2026-10-02
\bibitem{Weil52} A. Weil. Sur les ``formules explicites'' de la th\'eorie des nombres premiers. \emph{Communications du S\'eminaire Math\'ematique de l'Universit\'e de Lund}, tome suppl\'ementaire (1952), 252--265.
% verified: CrossRef 10.1103/physrev.96.1124 2026-10-09
\bibitem{Wick54} G. C. Wick. Properties of Bethe-Salpeter wave functions. \emph{Physical Review} 96 (1954), 1124--1134.
% verified: zbMATH 0063.08245 2026-09-27
\bibitem{Widder41} D. V. Widder. \emph{The Laplace Transform}. Princeton Mathematical Series 6, Princeton University Press, 1941.
% verified: CrossRef 10.2307/1968102 2026-09-26; pages 1 to 100 from zbMATH 0004.05905, 2026-10-02
\bibitem{Wiener32} N. Wiener. Tauberian theorems. \emph{Annals of Mathematics} 33 (1932), 1--100.
% verified: CrossRef 10.2307/1970008 2026-10-10 (CrossRef gives the first page only; pages 325 to 327 from zbMATH 0085.13203)
\bibitem{Wig58} E. P. Wigner. On the distribution of the roots of certain symmetric matrices. \emph{Annals of Mathematics} 67 (1958), 325--327.
% verified: zbMATH 0744.60100 2026-09-26
\bibitem{Williams90} D. Williams. Brownian motion and the Riemann zeta-function. In \emph{Disorder in Physical Systems}, G. R. Grimmett and D. J. A. Welsh, eds., Oxford University Press, 1990, 361--372.
% verified: CrossRef 10.1215/s0012-7094-56-02317-1 2026-09-26; pages 189 to 207 from zbMATH 0070.28501, 2026-10-02
\bibitem{Williamson} R. E. Williamson. Multiply monotone functions and their Laplace transforms. \emph{Duke Mathematical Journal} 23 (1956), 189--207.
% verified: CrossRef 10.1016/0304-4149(82)90050-3 2026-09-26; pages 301 to 312 from OpenAlex, 2026-10-02
\bibitem{Wolfe82} S. J. Wolfe. On a continuous analogue of the stochastic difference equation $X_n=\rho X_{n-1}+B_n$. \emph{Stochastic Processes and their Applications} 12 (1982), 301--312.
% verified: CrossRef 10.1214/aop/1176995474 2026-09-26; pages 523 to 531 from zbMATH 0394.60017, 2026-10-02
\bibitem{Yamazato78} M. Yamazato. Unimodality of infinitely divisible distribution functions of class $L$. \emph{Annals of Probability} 6 (1978), 523--531.
% verified: CrossRef 10.1007/bf00531612 2026-10-09
\bibitem{Yor80} M. Yor. Loi de l'indice du lacet brownien, et distribution de Hartman-Watson. \emph{Zeitschrift f\"ur Wahrscheinlichkeitstheorie und Verwandte Gebiete} 53 (1980), 71--95.
% verified: zbMATH 0817.11041 2026-09-27 (DOI 10.2969/aspm/02110281; editors from zbMATH 0771.00036; conference at Tokyo Institute of Technology, August 14 to 18, 1990, per Project Euclid)
\bibitem{Yoshida92} H. Yoshida. On Hermitian forms attached to zeta functions. In \emph{Zeta Functions in Geometry (Tokyo, 1990)}, N. Kurokawa and T. Sunada, eds., Advanced Studies in Pure Mathematics 21, Kinokuniya, 1992, 281--325.
\end{thebibliography}
